% P06.3 literal harmonic-approximant counterexample, exact source-aware
% packet, 25 August 2026.  No novelty or priority claim is intended.

\paragraph{Source scope and distinction of the three classes.}
The proof uses the following precise source statements.
\begin{enumerate}
\item Chambert--Loir--Ducros smooth functions are ordinary smooth
$(0,0)$-forms: locally they are pullbacks of smooth Lagerberg functions by
moment maps whose coordinates are logarithms of analytic units.  This is
the definition used in Chambert--Loir--Ducros, arXiv:1204.6277v2,
\S\S3.1.7 and 3.4.7, and in Wanner, arXiv:1801.04713v3,
Definitions~3.9--3.12.
\item Gubler--Jell--Rabinoff, arXiv:2111.05741v3,
Proposition~10.14 embeds ordinary CLD-smooth forms into weakly smooth forms.
Proposition~10.17 says that weakly smooth functions of zero
$\mathrm d'\mathrm d''$ are precisely their $\mathbf R$-harmonic
functions, and Proposition~15.7 identifies this harmonicity with
Thuillier's on rig-smooth curves.  Remark~15.10 identifies the ordinary
CLD-smooth harmonic
subsheaf $\mathcal F$ with the $\mathbf R$-span of germs $\log|f|$ for
analytic units.  Thuillier, Lemma~2.3.22, as recalled and used in the proof
of Gubler--Jell--Rabinoff, Proposition~17.4, gives
\[
 (\mathcal H/\mathcal F)_P
       \simeq \operatorname{Pic}^0(E)\otimes_{\mathbf Z}\mathbf R
\]
at a type--$2$ point with residual curve $E$.
Proposition~15.11 regularizes an arbitrary subharmonic function on a
relatively compact open subset of any rig-smooth boundaryless curve by a
decreasing net of \emph{weakly smooth psh} functions.  It does not say that
the approximants are ordinary CLD-smooth or harmonic.
\item Gubler--Jell--Rabinoff, \S17, constructs the smooth proper genus-$3$
curve and the harmonic function used below.  In the notation of their
Proposition~17.4, the residual curve at the vertex $P$ is an elliptic curve
$E$ with identity $z$ and a non-torsion point $x$, and one coordinate
$\widetilde T_1$ of the lifted tropical Abel--Jacobi map is integral PL and
harmonic near $P$.  Its residual slope divisor represents
$[z]-[x]\in\operatorname{Pic}^0(E)\otimes\mathbf R$, so it is not
ordinary CLD-smooth.
\item Wanner, Corollary~5.4, proves local decreasing approximation by
ordinary CLD-smooth \emph{psh} functions when the residue field is
algebraic over a finite field, or the curve is the projective line or a
Mumford curve.  Remark~5.5 says that there are curves with
$\mathcal F\subsetneq\mathcal H$ and that ordinary CLD-smooth psh
approximability is not known there.  Hence non-smoothness of a harmonic
function alone is not a counterexample to psh approximation; the closed
finite-support argument below is needed for the literal harmonic target.
\end{enumerate}
The audited arXiv source archives for Wanner v3 and Gubler--Jell--Rabinoff
v3 have SHA--256 values
\texttt{041d4032e4dba67dc8677981a0f40ac0230ad3404460ea05bd22f7e86d641686}
and
\texttt{58f64063539ade89ada2be28a7b9e4ceae7f0b2beb3786d70cac062645a31789},
respectively.  The previously audited CLD v2 PDF has SHA--256
\texttt{5690f999523017d184ab751bb55857ab1637b197d1186ad4fd2edba7e1bda719}.

\begin{proposition}
There exist a complete algebraically closed nontrivially valued
non-Archimedean field $K$, a smooth proper connected $K$-curve $C$, a
type--$2$ point $P\in C^{\mathrm{an}}$, a relatively compact connected open
neighbourhood $U$ of $P$, and an integral PL harmonic function
$h:U\to\mathbf R$ with the following properties.
\begin{enumerate}
\item The function $h$ is weakly smooth and
\[
                         \mathrm d'\mathrm d''h=0
\]
as a weakly smooth form.
\item There is no net $(h_i)$ of ordinary Chambert--Loir--Ducros-smooth
harmonic functions on $U$---that is,
$\mathrm d'\mathrm d''h_i=0$ as an ordinary smooth form---which converges
locally uniformly to $h$.
In particular, there is no decreasing net of ordinary
Chambert--Loir--Ducros-smooth harmonic functions on $U$ which converges
pointwise to $h$.
\end{enumerate}
Thus the literal harmonic-approximant reading of the third question in
\textup{\texttt{prob:frontier-pluriharmonic}} is false, already on a smooth
proper Berkovich curve.  This does not disprove decreasing regularization by
ordinary Chambert--Loir--Ducros-smooth psh functions: those approximants are
not required to have zero curvature, and that stronger regularization
problem remains outside the proposition.
\end{proposition}

\begin{proof}
We use the curve in Gubler--Jell--Rabinoff, \S17.  The ground field $K$ is
complete, algebraically closed and nontrivially valued, with algebraically
closed residue field not algebraic over a finite field.  Their metrized
complex has a theta graph with a vertex $P$.  The residual curve at $P$ is
an elliptic curve $E$; the three graph directions $a,b,c$ at $P$ correspond
to three distinct points $x,y,z\in E(\widetilde K)$, where $z$ is the
identity and $x$ is non-torsion.  The realization theorem cited in their
construction gives a smooth proper connected genus-$3$ curve $C/K$ with
this metrized complex.

Let $h=\widetilde T_1$ be the first coordinate in their
Proposition~17.4.  After shrinking its domain, it is an integral PL
harmonic function near $P$.  Its residual slope divisor $D_P(h)$ is
supported on
\[
                         S=\{x,y,z\}\subset E(\widetilde K)
\tag{1}\label{eq:p063-support}
\]
and its class is
\[
 [D_P(h)]=[z]-[x]
   \ne0\quad\text{in}\quad
 \operatorname{Pic}^0(E)\otimes_{\mathbf Z}\mathbf R.
\tag{2}\label{eq:p063-nonzero-class}
\]
The last inequality follows from the choice of the non-torsion point $x$.
Indeed, the class $[z]-[x]$ then has infinite order; the injection of its
infinite cyclic subgroup into $\operatorname{Pic}^0(E)$ remains injective
after tensoring with $\mathbf R$, because $\mathbf R$ is a torsion-free,
hence flat, $\mathbf Z$-module.
By Gubler--Jell--Rabinoff, Proposition~10.17, $h$ is weakly smooth and
$\mathrm d'\mathrm d''h=0$.

Choose points $Q_a,Q_b,Q_c$ in the interiors of the three graph edges,
sufficiently close to $P$ that the closed three-armed star domain $Y_0$ cut
out by them lies in the domain on which $h$ is harmonic.  Its analytic
boundary is $\{Q_a,Q_b,Q_c\}$.  Put
\[
                         U=Y_0\setminus\partial Y_0.
\tag{3}\label{eq:p063-open-star}
\]
After moving the $Q_e$ slightly toward $P$ if necessary, $Y_0$ is a compact
strictly affinoid domain and $U$ is a relatively compact connected open
neighbourhood of $P$.

Choose points $R_e\in(P,Q_e)$ and let $Y\Subset U$ be the compact inner
star with endpoints $R_a,R_b,R_c$.  We record a fixed-support consequence
of harmonicity.  Subdividing the stable model at $R_a,R_b,R_c$ and taking
the corresponding formal open gives a strictly semistable formal model
$\mathfrak Y$ of $Y$ whose skeleton is exactly the three segments
$[P,R_e]$ and whose boundary is $\{R_a,R_b,R_c\}$.  Let $g$ be harmonic on
$U$ in the sense of Gubler--Jell--Rabinoff.  By their Proposition~15.7,
$g|_Y$ belongs to Thuillier's space $H(Y)$.  The independence of
$H(Y)$ from the chosen strictly semistable model (Thuillier,
Proposition~2.3.3, recalled in Wanner, \S2) therefore gives
$g|_Y=G\circ\tau_{\mathfrak Y}$, where $G$ is piecewise affine on
$S(\mathfrak Y)$ and has zero sum of outgoing slopes away from
$\partial Y$.  At every subdivision point in the interior of an arm, this
equation says that the two one-sided slopes agree.  Hence $G$ is affine
with constant slope on each $[P,R_e]$.  Since $g$ factors through this
three-armed retraction, all other tangent directions at $P$ have slope
zero.  Consequently the slope divisor of $g$ at $P$ is supported on the
three tangent directions corresponding to $S$, and
\[
       D_P(g)\in
       V_S:=\operatorname{Div}^0_S(E)\otimes_{\mathbf Z}\mathbf R,
\tag{4}\label{eq:p063-fixed-support-space}
\]
and the three coefficients of $D_P(g)$ are the outgoing slopes of $g$ on
the fixed arms of $Y$.  Here $D_P(g)$ denotes the divisor on the residual
curve whose coefficients are these outgoing slopes; its class is the
residue class of the germ of $g$ at $P$.

If, in addition, $g$ is ordinary CLD-smooth, then Gubler--Jell--Rabinoff,
Remark~15.10, gives $g_P\in\mathcal F_P$.  Thus its residual divisor is a
real linear combination of principal divisors, or equivalently its class
in $\operatorname{Pic}^0(E)\otimes_{\mathbf Z}\mathbf R$ vanishes.
Combining this with
\eqref{eq:p063-fixed-support-space} gives
\[
 D_P(g)\in
 P_S:=\ker\!\left(
 V_S\longrightarrow\operatorname{Pic}^0(E)\otimes_{\mathbf Z}\mathbf R
 \right).
\tag{5}\label{eq:p063-principal-subspace}
\]
The space $V_S$ is finite-dimensional and $P_S$ is a linear subspace;
hence $P_S$ is closed in the coefficient topology of $V_S$.

Suppose first that a net $(h_i)$ of ordinary CLD-smooth harmonic functions
on $U$ converges locally uniformly to $h$; here harmonic means
$\mathrm d'\mathrm d''h_i=0$ as an ordinary smooth form.  By the
inclusion of ordinary into weakly smooth forms and
Gubler--Jell--Rabinoff, Proposition~10.17, every $h_i$ is harmonic in their
sense.  By the fixed-support argument,
$h_i$ and $h$ are
affine on each $[P,R_e]$.  If $\ell_e$ is its length, then the outgoing
slope is
\[
          \frac{h_i(R_e)-h_i(P)}{\ell_e},
\tag{6}\label{eq:p063-slope-from-values}
\]
and the analogous formula holds for $h$.  Uniform convergence on $Y$
therefore gives coefficientwise convergence
\[
                         D_P(h_i)\longrightarrow D_P(h)
                         \quad\text{in }V_S.
\tag{7}\label{eq:p063-divisor-convergence}
\]
Every term on the left lies in the closed subspace $P_S$ by
\eqref{eq:p063-principal-subspace}.  Hence $D_P(h)\in P_S$, contrary to
\eqref{eq:p063-nonzero-class}.

Finally, suppose that $(h_i)$ is a decreasing net of ordinary CLD-smooth
harmonic functions on $U$ with pointwise limit $h$.  On the compact set
$Y$, the net version of Dini's theorem makes the convergence uniform.  To
see this directly, for $\varepsilon>0$ the increasing family of open sets
\[
              \{p\in Y:h_i(p)-h(p)<\varepsilon\}
\]
covers $Y$; compactness supplies finitely many indices and directedness
supplies one index dominating all of them.  The preceding contradiction
then applies.  This proves both assertions.
\end{proof}

\paragraph{Scope firewall.}
The finite-dimensional closure argument uses zero curvature of every
approximant.  It does not apply to ordinary CLD-smooth psh approximants:
their positive curvature may move and concentrate, and uniform convergence
does not force endpoint derivatives of convex functions to converge.
Accordingly the proposition neither contradicts Gubler--Jell--Rabinoff,
Proposition~15.11, nor resolves the general CLD-smooth psh approximation
question left open in Wanner, Remark~5.5.  No novelty or priority claim is
made.
